% Study sheet: Volume IV, Chapter 3, The Cognitive Brain.
% Numbers: \val{...} from brain_studysheet_numbers.tex, written by brain_studysheet_numbers.py.
\documentclass{studysheet}
\input{brain_studysheet_numbers}
\sheetvolume{4}
\sheetchapter{3}{The Cognitive Brain}
\sheetsources{Numbers from the chapter's mlsysim scenarios for the warehouse mobile manipulator; questions marked assumed use new inputs.}

\begin{document}
% ===================================================================== FRONT: LEARN
\frontside
\begin{bigidea}
A learned model handles objects no engineer listed, but it is the slowest computation aboard
and only \hl{proposes}. Two budgets govern it: how often it \hl{renews} and how long each answer
stays \hl{true}. Motion cannot be recalled, so every proposal stays refusable.
\end{bigidea}

\begin{leftcol}
\sheetsec{Key ideas}
\begin{ideas}
\idea{Learned proposes; the permission path admits.}{Classical geometry is provable but runs out of cases;
  a learned model covers more, proves nothing.}
\idea{Prefill and decode bind differently.}{Prefill reuses each weight across the whole input
  (compute-bound); decode rereads all weights per token.}
\idea{Chunks supply motion, not fresh evidence.}{$\ActionChunkHorizon$ setpoints per evaluation cut weight traffic
  per setpoint; new evidence still waits a renewal plus inference.}
\idea{Two proposers.}{A slow intent model names a target pose; a fast chunk policy renews motion. What
  crosses between them sets what can be checked.}
\idea{On time can still be wrong.}{An in-limits proposal can aim at the wrong
  thing; only outside evidence refuses it.}
\end{ideas}

\vspace{9pt}
\fcard{Memory wall}%
  {\ProposerEvaluationRate \le \frac{\SustainedMemoryBandwidth}{\WeightMemory}, \qquad \WeightMemory = \Parameters \times \text{bytes per parameter}}%
  {At batch size one each evaluation streams all $\Parameters$ weights: bytes moved, not arithmetic, set the rate.}%
  {}{}
\fcard{One fresh evaluation}%
  {\ProposerInferenceTime \approx \under{\max\!\Big(\frac{\PrefillOperations}{\EfficiencyHw\PeakRate},\ \frac{\WeightMemory}{\SustainedMemoryBandwidth}\Big)}{prefill}
   + \under{(\EmittedTokenCount - 1)\,\frac{\WeightMemory}{\SustainedMemoryBandwidth}}{decode sweeps}}%
  {$\PrefillOperations \approx 2\Parameters\SequenceLength$ ($\SequenceLength$ input tokens, efficiency $\EfficiencyHw$
   of peak $\PeakRate$, $\EmittedTokenCount$ output tokens). Weight-only quantization shrinks only decode.}%
  {Prefill intensity is about $\SequenceLength$ FLOP/byte at FP16, decode about one; above the ridge
   ($\PeakRate$ over peak bandwidth), compute binds.}{}
\end{leftcol}%
\begin{rightcol}
\sheetsec{A chunk is buffer, not permission}
{\centering
\begin{tikzpicture}[sheet]
  \pgfmathsetmacro\s{236/\val{span}}
  \pgfmathsetmacro\xa{\val{infer}*\s}\pgfmathsetmacro\xb{\val{age}*\s}
  \pgfmathsetmacro\xc{\val{tau}*\s}\pgfmathsetmacro\xd{\val{span}*\s}
  \fill[rulegray] (0,20) rectangle (\xa,32);
  \draw[hotbox] (\xa,20) rectangle (\xb,32);
  \draw[accbox] (\xb,20) rectangle (\xc,32);
  \draw[nbox, dashed] (\xc,20) rectangle (\xd,32);
  \draw[dim, densely dashed] (\xc,18) -- (\xc,50);
  \node[lab, anchor=south] at ({(\xa+\xb)/2},34) {in force};
  \node[lab, anchor=south] at ({(\xb+\xc)/2},34) {reserve if leased};
  \draw[dim] (0,50) -- (0,53) -- (\xc,53) -- (\xc,50);
  \node[lab, anchor=south] at ({\xc/2},54) {evidence holds, capture to $\EvidenceHorizon$};
  \node[lab, anchor=south] at ({(\xc+\xd)/2},34) {buffer, not permission};
  \node[small, anchor=north] at (0,17) {0};
  \node[small, anchor=south] at ({\xa/2},34) {infer};
  \node[small, anchor=north] at (\xa,17) {\val{infer}};
  \node[small, anchor=north] at (\xb,17) {\val{age}};
  \node[small, anchor=north] at (\xc,12) {$\EvidenceHorizon$ = \val{tau}};
  \node[small, anchor=north] at (\xd,17) {\val{span} ms};
\end{tikzpicture}\par}
\vspace{3pt}
{\fontsize{8.6}{11}\selectfont From capture: $\ActionChunkHorizon=\val{horizon}$ setpoints \val{step} ms apart, renewed every \val{renew} ms.\par}

\vspace{6pt}
\fcard{Evidence horizon and minimum supply}%
  {\EvidenceHorizon = \sqrt{\frac{2(\PositionalTaskTolerance - \InitialTrackingError)}{\SlipAccelerationBound}},
   \qquad \MinimumChunkSetpointCount = \left\lceil \frac{\ChunkRenewalPeriod + \ProposerInferenceTime}{\SetpointStepInterval} \right\rceil}%
  {Slip $\SlipAccelerationBound$ after capture eats the margin from initial
   error $\InitialTrackingError$ to tolerance $\PositionalTaskTolerance$. Need $\ChunkRenewalPeriod + \ProposerInferenceTime \le \EvidenceHorizon$.}%
  {$\MinimumChunkSetpointCount$ steps of $\SetpointStepInterval$ bridge one renewal $\ChunkRenewalPeriod$.}{}
\fcard{KV cache of retained context}%
  {\KvCacheMemory = 2 \times \LayerCount \times \HiddenDimension \times \SequenceLength \times \ElementStorageSize}%
  {$\SequenceLength$ tokens kept, $\LayerCount$ layers of width $\HiddenDimension$, $\ElementStorageSize$ bytes per value
   (two at FP16). One read takes $\KvCacheMemory/\SustainedMemoryBandwidth$.}{}{}

\vspace{6pt}
\sheetsec{Decoders: one fresh \val{horizon}-step chunk}
\begin{sheettable}{@{}l L r@{}}
\thead{Decoder} & \thead{After prefill} & \thead{Total}\\ \midrule
Autoregressive & \val{decsweeps} sequential sweeps & \val{decbinnedsec} s\\
One-pass head & none (longer prefill) & \val{decparallel} ms\\
Iterative expert & \val{decsteps} small-expert steps & \val{decexpert} ms\\
\end{sheettable}
\vspace{3pt}
{\fontsize{8.6}{11}\selectfont Each first pays the \val{tpre} ms prefill.\par}

\vspace{8pt}
\sheetsec{Refuse on outside evidence}
\begin{sheettable}{@{}>{\raggedright\arraybackslash}p{0.40\linewidth} L@{}}
\thead{What arrives} & \thead{What can refuse it}\\ \midrule
Latent vector & Only its motion; no position or age\\
Target pose + capture time & Depth; its age against $\EvidenceHorizon$\\
Grasp on a printed handle & Depth shows a flat surface\\
Confidence near one & Depth vetoes; softmax sums to one by design\\
Label flipping (chattering) & Depth and track hold still; debounce (ages evidence), filter, limit jerk\\
\end{sheettable}
\end{rightcol}

% ===================================================================== BACK: PRACTICE
\backside
\begin{answers}
\ans{Finish}{\val{tpre} + \val{decodesteps} $\times$ \val{sweepeight} = \val{tpre} + \val{deceight} = \res{\val{toteight} ms}:
  fits, \val{slackeight} ms spare. Faster needs fewer bytes or a shorter prefill; shrinking weights alone cannot
  pass \val{ceiling} Hz, the prefill limit (evaluation card).}
\ans{Q1}{(a) $\sqrt{\val{qtauinner}}$ s $\approx$ \val{qtau} ms. (b) \val{renew} + \val{infer} = \val{age} ms fits;
  $\MinimumChunkSetpointCount = \lceil \val{age}/\val{step} \rceil = \val{hmin}$. (c) \val{qrenewslow} + \val{infer} =
  \val{qageslow} ms $>$ \val{qtau}: evidence expires first (evidence card).}
\ans{Q2}{\val{pertoken} B $\times$ \val{ctx} $\approx$ \val{perstep} MB per step, $\times$ \val{qsteps} = \val{qcache} GB;
  with weights \val{qtotal} GB, \val{qshare}\% of \val{memgb} GB. One full read takes \val{qcache} $\div$ \val{bw}
  $\approx$ \val{qreread} ms of shared bandwidth, longer than a whole \val{renew} ms chunk renewal (KV card).}
\columnbreak
\ans{Q3}{The target: a latent holds no position to compare with depth or age against drift; only chunk
  motion stays checkable. Keep a target pose with capture time (table).}
\ans{Q4}{Perceptual chattering: the label flips while depth and tracked position hold
  still. Debounce, filter, limit jerk; debouncing delays real changes and ages the evidence (table).}
\end{answers}

\begin{leftcol}
\sheetsec{Worked example}
\begin{practicebox}\boxtitle{How often can the intent model renew?}
The running machine's intent model, \val{params} billion parameters: bandwidth \val{bwpeak} GB/s peak, \val{buspct}\% sustained
(\val{bw} GB/s); \val{peak} TFLOP/s at FP16, $\EfficiencyHw$ = \val{util}; $\SequenceLength$ = \val{tokensin},
$\PrefillOperations$ = \val{flops}$\times10^{12}$ FLOPs (given), $\EmittedTokenCount$ = \val{tokensout}.\par\vspace{5pt}
{\centering
\begin{tikzpicture}[sheet]
  \pgfmathsetmacro\sc{150/\val{totfp}}
  \pgfmathsetmacro\p{\val{tpre}*\sc}
  \foreach \lab/\d/\tot/\y in {FP16/\val{decfp}/\val{totfp}/16, INT4/\val{decint}/\val{totint}/0}{
    \node[lab, anchor=east] at (30,\y+6) {\lab};
    \draw[accbox] (34,\y) rectangle ({34+\p},\y+12);
    \draw[nbox] ({34+\p},\y) rectangle ({34+\tot*\sc},\y+12);
    \node[note, anchor=west] at ({38+\tot*\sc},\y+6) {\tot\ ms};}
  \node[lab] at ({34+\p/2},38) {prefill};
  \node[lab] at ({34+\p+\val{decint}*\sc/2+10},38) {decode};
  \draw[dim, densely dashed] ({34+\val{period}*\sc},-4) -- ({34+\val{period}*\sc},32);
  \node[small, anchor=north] at ({34+\val{period}*\sc},-8) {\val{period} ms period};
\end{tikzpicture}\par}\vspace{4pt}
\begin{steps}
\item FP16 weights: \val{params} billion $\times$ 2 bytes = \val{wfp} GB; sweep \val{wfp} $\div$ \val{bw} = \val{sweepfp} ms, a wall of \val{wallfp} Hz.
\item Prefill: \val{tokensin} FLOP/byte $>$ ridge \val{peak}T $\div$ \val{bwpeak}G $\approx$ \val{ridge}, so compute-bound:
  \val{flops}$\times10^{12} \div (\val{util} \times \val{peak}\times10^{12})$ = \res{\val{tpre} ms}.
\item FP16: \val{tpre} + \val{decodesteps} $\times$ \val{sweepfp} = \val{tpre} + \val{decfp} = \res{\val{totfp} ms}.
\item INT4 (\val{wint} GB, sweep \val{sweepint} ms): \val{tpre} + \val{decodesteps} $\times$ \val{sweepint} = \res{\val{totint} ms},
  so the machine runs at \val{rate} Hz.
\end{steps}
\takeaway{A quarter of the bytes gives only \val{speedup}$\times$ (\val{lesstime}\% less time): prefill is now \val{preshare}\%.}
\end{practicebox}

\vspace{10pt}
\sheetsec{Your turn: finish it}
\begin{practicebox}\boxtitle{Does an 8-bit build fit the same period?}
One byte per weight.\par\vspace{3pt}
\textbf{Given:} \val{weight} GB; sweep \val{weight} $\div$ \val{bw} = \val{sweepeight} ms.\par
\textbf{You:} find the evaluation time; does it fit \val{period} ms? What would a faster rate need?\par\vspace{30pt}
\end{practicebox}
\end{leftcol}%
\begin{rightcol}
\sheetsec{Common mistakes}
\mistake{``A smooth chunk inside joint limits is safe to run to its end.''}%
  {Its \val{span} ms rest on one capture whose evidence holds \val{tau} ms. Run to the end, \val{slip} m/s$^2$
   slip drifts \val{driftend} mm: \val{errend} mm of error against \val{tol} mm tolerance.}{}
\mistake{``With chunks and a fast decoder, the big model can drive the arm.''}%
  {Every decoder first pays the \val{tpre} ms prefill (decoder table), above the
   \val{renew} ms renewal. A fresh chunk lands \val{totint} ms after capture, \val{gap} renewal periods.}{}

\mistake{``We must choose: provable classical control or a capable learned policy.''}%
  {Not either-or: the learned policy proposes; an independent filter corrects any command that would cross a
   limit.}{}

\vspace{2pt}
\sheetsec{Check yourself}
\begin{qs}
\item \assumed{} slip \val{qslip} m/s$^2$; tolerance \val{tol} mm, initial error \val{errzero} mm, \val{step} ms steps,
  \val{infer} ms inference. (a) $\EvidenceHorizon$? (b) Does \val{renew} ms renewal fit; $\MinimumChunkSetpointCount$? (c) Renewal slowed to \val{qrateslow} Hz?
\item \assumed{} a \val{qwindow} s history at \val{qhz} Hz: \val{layers} layers, $\HiddenDimension$ = \val{width}, FP16,
  \val{ctx} tokens per step. Cache size; share of \val{memgib} GiB (\val{memgb} GB) with the \val{wint} GB INT4 weights? Read time at \val{bw} GB/s?
\item A reviewer proposes passing the chunk policy a latent vector instead of a target pose: richer, no
  decoding. What does the permission path lose?
\item Motor currents spike; the log shows the target's label flipping between ``mug'' and ``bowl'' every
  frame. Name the failure, the confirming evidence, and the defenses with their cost.
\end{qs}
\end{rightcol}
\end{document}
